% =====================================================================
% ASL AI Education Platform -- B.Tech Major Project Report
% SRM Institute of Science and Technology, Kattankulathur
% Authors: Jasmine (RA2211003010059), Arihant Jain (RA2211003010105)
% Submission: April 2026
% =====================================================================

\documentclass[12pt,a4paper,oneside]{report}

\usepackage[a4paper,left=1.25in,right=1in,top=1in,bottom=1in]{geometry}
\usepackage{times}
\usepackage{setspace}
\usepackage{graphicx}
\usepackage{tabularx}
\usepackage{array}
\usepackage{longtable}
\usepackage{booktabs}
\usepackage{xcolor}
\usepackage{titlesec}
\usepackage{tocloft}
\usepackage{fancyhdr}
\usepackage{url}
\usepackage{hyperref}
\usepackage{enumitem}
\usepackage{indentfirst}
\usepackage{caption}
\usepackage{listings}
\usepackage{float}

\onehalfspacing
\setlength{\parindent}{1.5em}
\setlength{\parskip}{0.4em}

\titleformat{\chapter}[block]
  {\normalfont\bfseries\Large\centering}
  {CHAPTER \thechapter}{1em}{\\\Large}
\titlespacing*{\chapter}{0pt}{-30pt}{30pt}

\titleformat{\section}{\normalfont\bfseries\large}{\thesection}{1em}{}
\titleformat{\subsection}{\normalfont\bfseries\normalsize}{\thesubsection}{1em}{}

\setcounter{tocdepth}{2}
\renewcommand{\cftchapdotsep}{\cftdotsep}
\renewcommand{\cftchapleader}{\cftdotfill{\cftchapdotsep}}

\pagestyle{fancy}
\fancyhf{}
\fancyfoot[C]{\thepage}
\renewcommand{\headrulewidth}{0pt}
\renewcommand{\footrulewidth}{0pt}
\fancypagestyle{plain}{%
  \fancyhf{}%
  \fancyfoot[C]{\thepage}%
  \renewcommand{\headrulewidth}{0pt}%
  \renewcommand{\footrulewidth}{0pt}%
}

\lstset{
  basicstyle=\ttfamily\small,
  breaklines=true,
  frame=single,
  captionpos=b,
  numbers=left,
  numberstyle=\tiny,
  keywordstyle=\bfseries,
  showstringspaces=false
}

\begin{document}

% =====================================================================
% TITLE PAGE
% =====================================================================
\begin{titlepage}
\thispagestyle{empty}
\begin{center}
\vspace*{0.2cm}
{\bfseries\large An ASL-Based Storytelling and AI-Powered Learning Platform for Hearing-Impaired Children Using Real Sign Language Video Data and Generative AI\par}

\vspace{0.8cm}
{\large A PROJECT REPORT\par}
\vspace{0.3cm}
{\large Submitted by\par}
\vspace{0.25cm}
{\bfseries\large JASMINE \quad [RA2211003010059]\par}
\vspace{0.15cm}
{\bfseries\large ARIHANT JAIN \quad [RA2211003010105]\par}

\vspace{0.6cm}
{\large \textit{Under the Guidance of}\par}
\vspace{0.2cm}
{\bfseries\large Dr.\ P.\ Rajasekaran\par}
{\normalsize (Assistant Professor, Department of Computing Technologies)\par}

\vspace{0.55cm}
{\large in partial fulfilment of the requirements for the degree of\par}
\vspace{0.25cm}
{\bfseries\large BACHELOR OF TECHNOLOGY\par}
{\large in\par}
{\bfseries\large COMPUTER SCIENCE AND ENGINEERING\par}

\vspace{0.6cm}
\includegraphics[width=0.20\textwidth]{srm_logo.png}\\[0.3cm]
{\bfseries DEPARTMENT OF COMPUTING TECHNOLOGIES\\
COLLEGE OF ENGINEERING AND TECHNOLOGY\\
SRM INSTITUTE OF SCIENCE AND TECHNOLOGY\\
KATTANKULATHUR -- 603\,203\par}

\vspace{0.5cm}
{\bfseries\large APRIL 2026\par}
\end{center}
\end{titlepage}
\setcounter{page}{1}

% =====================================================================
% OWN WORK DECLARATION
% =====================================================================
\thispagestyle{empty}
\begin{center}
{\bfseries Department of Computing Technologies\\
SRM Institute of Science \& Technology\\[0.4em]
Own Work Declaration Form\par}
\end{center}

\noindent\small
This sheet must be filled in (each box ticked to show that the condition has been met). It must be signed and dated along with your student registration number and included with all assignments you submit -- work will not be marked unless this is done.

\vspace{0.4em}
\noindent\textbf{To be completed by the student for all assessments}

\vspace{0.6em}
\noindent\begin{tabularx}{\linewidth}{@{}lX@{}}
\textbf{Degree / Course:} & B.Tech in Computer Science and Engineering \\[0.3em]
\textbf{Student Names:} & Jasmine, Arihant Jain \\[0.3em]
\textbf{Registration Numbers:} & RA2211003010059, RA2211003010105 \\[0.3em]
\textbf{Title of Work:} & An ASL-Based Storytelling and AI-Powered Learning Platform for Hearing-Impaired Children Using Real Sign Language Video Data and Generative AI \\
\end{tabularx}

\vspace{0.8em}
\noindent We hereby certify that this assessment complies with the University's Rules and Regulations relating to Academic Misconduct and Plagiarism, as listed in the University Website, Regulations, and the Education Committee guidelines.

\vspace{0.4em}
\noindent We confirm that all the work contained in this assessment is our own except where indicated, and that we have met the following conditions:

\begin{itemize}[leftmargin=2em,itemsep=0.1em,topsep=0.2em]
  \item[$\square$] Clearly referenced / listed all sources as appropriate.
  \item[$\square$] Referenced and put in inverted commas all quoted text (from books, web, etc.).
  \item[$\square$] Given the sources of all pictures, data etc.\ that are not our own.
  \item[$\square$] Not made any use of the report(s) or essay(s) of any other student(s) either past or present.
  \item[$\square$] Acknowledged in appropriate places any help that we have received from others.
  \item[$\square$] Complied with any other plagiarism criteria specified in the Course handbook / University website.
\end{itemize}

\noindent We understand that any false claim for this work will be penalised in accordance with the University policies and regulations.

\vspace{0.4em}
\noindent\textbf{DECLARATION:} We are aware of and understand the University's policy on Academic Misconduct and Plagiarism, and we certify that this assessment is our own work, except where indicated by referencing, and that we have followed the good academic practices noted above.

\vspace{1.6em}
\noindent\begin{tabularx}{\linewidth}{@{}lXr@{}}
RA2211003010059 \quad : \quad \rule{6cm}{0.4pt} & & Date: \rule{2.5cm}{0.4pt} \\[1.2em]
RA2211003010105 \quad : \quad \rule{6cm}{0.4pt} & & Date: \rule{2.5cm}{0.4pt} \\
\end{tabularx}

\normalsize
\clearpage

% =====================================================================
% BONAFIDE CERTIFICATE
% =====================================================================
\thispagestyle{empty}
\begin{center}
{\bfseries SRM INSTITUTE OF SCIENCE AND TECHNOLOGY\\
KATTANKULATHUR -- 603\,203\par}
\vspace{0.8cm}
{\bfseries\large BONAFIDE CERTIFICATE\par}
\end{center}

\vspace{0.6cm}
\noindent Certified that 21CSP302L\,--\,Project Report titled \textbf{``An ASL-Based Storytelling and AI-Powered Learning Platform for Hearing-Impaired Children Using Real Sign Language Video Data and Generative AI''} is the bonafide work of \textbf{``JASMINE [RA2211003010059]\,, ARIHANT JAIN [RA2211003010105]''} who carried out the project work under my supervision. Certified further that, to the best of my knowledge, the work reported herein does not form any other project report or dissertation on the basis of which a degree or award was conferred on an earlier occasion on this or any other candidate.

\vspace{2.5cm}
\noindent\begin{tabularx}{\linewidth}{@{}XX@{}}
\hfill\rule{5cm}{0.4pt}\hfill & \hfill\rule{5cm}{0.4pt}\hfill \\[0.2em]
\hfill SIGNATURE\hfill & \hfill SIGNATURE\hfill \\[1.2em]
\hfill\textbf{Dr.\ P.\ Rajasekaran}\hfill & \hfill\textbf{Dr.\ G.\ Niranjana}\hfill \\
\hfill SUPERVISOR\hfill & \hfill PROFESSOR \& HEAD\hfill \\
\hfill Assistant Professor\hfill & \hfill DEPARTMENT OF\hfill \\
\hfill Department of Computing Technologies\hfill & \hfill COMPUTING TECHNOLOGIES\hfill \\
\end{tabularx}

\vspace{2.5cm}
\noindent\begin{tabularx}{\linewidth}{@{}XX@{}}
\hfill\rule{5cm}{0.4pt}\hfill & \hfill\rule{5cm}{0.4pt}\hfill \\[0.2em]
\hfill EXAMINER 1\hfill & \hfill EXAMINER 2\hfill \\
\end{tabularx}
\clearpage

% =====================================================================
% ACKNOWLEDGEMENTS
% =====================================================================
\pagenumbering{roman}
\setcounter{page}{1}

\chapter*{ACKNOWLEDGEMENTS}
\addcontentsline{toc}{chapter}{ACKNOWLEDGEMENTS}

\begin{singlespace}
We express our humble gratitude to \textbf{Dr.\ C.\ Muthamizhchelvan}, Vice-Chancellor, SRM Institute of Science and Technology, for the facilities extended for the project work and his continued support.

We extend our sincere thanks to \textbf{Dr.\ Leenus Jesu Martin M}, Dean-CET, SRM Institute of Science and Technology, for his invaluable support.

We wish to thank \textbf{Dr.\ Revathi Venkataraman}, Professor and Chairperson, School of Computing, SRM Institute of Science and Technology, for her support throughout the project work.

We encompass our sincere thanks to \textbf{Dr.\ M.\ Pushpalatha}, Professor and Associate Chairperson -- CS, School of Computing, and \textbf{Dr.\ Lakshmi}, Professor and Associate Chairperson -- AI, School of Computing, SRM Institute of Science and Technology, for their invaluable support.

We are incredibly grateful to our Head of the Department, \textbf{Dr.\ G.\ Niranjana}, Professor and Head, Department of Computing Technologies, SRM Institute of Science and Technology, for her suggestions and encouragement at all the stages of the project work.

We thank our Project Coordinators, Panel Head, Panel Members, and Faculty Advisors \textbf{Dr.\ Vanusha D} and \textbf{Dr.\ B.\ Arthi}, Department of Computing Technologies, SRM Institute of Science and Technology, for their inputs during the project reviews and continued support throughout our course.

Our inexpressible respect and thanks to our guide, \textbf{Dr.\ P.\ Rajasekaran}, Assistant Professor, Department of Computing Technologies, SRM Institute of Science and Technology, for providing us with the opportunity to pursue our project under his mentorship and for the freedom to explore the research topics of our interest.

We sincerely thank all the staff members of the Department of Computing Technologies, School of Computing, SRM Institute of Science and Technology, for their help during our project. Finally, we thank our parents, family members, and friends for their unconditional love, constant support, and encouragement.

\vspace{1.5em}
\begin{flushright}
\begin{tabular}{r}
\textbf{Jasmine} \quad [RA2211003010059]\\[0.4em]
\textbf{Arihant Jain} \quad [RA2211003010105]
\end{tabular}
\end{flushright}
\end{singlespace}

% =====================================================================
% ABSTRACT
% =====================================================================
\chapter*{ABSTRACT}
\addcontentsline{toc}{chapter}{ABSTRACT}

Access to inclusive and accessible education remains a significant challenge for hearing-impaired children, primarily due to the dominance of text-based and audio-centric learning resources. This report describes an \textbf{ASL-Based Storytelling and AI-Powered Learning Platform} that converts educational lesson content into both real American Sign Language (ASL) video narratives and AI-generated visual concept cards, making structured learning interactive and accessible for deaf and hard-of-hearing children.

The system combines two complementary pipelines: a \textit{real-sign video pipeline} that retrieves and concatenates authentic ASL clips from the Word-Level American Sign Language (WLASL) dataset using vocabulary matching and FFmpeg-based video processing, and an \textit{AI lesson generation pipeline} powered by Google Gemini 2.5 Flash that simplifies lesson text into ASL-friendly word sequences and triggers contextual image generation via the Stability AI SDXL 1.0 model or a locally rendered Pillow-based fallback.

The platform is delivered as a React 18 single-page application backed by a FastAPI service, with user data persisted in SQLite and all generated media served as static files. The two AI integrations --- Gemini for natural language simplification and Stability AI for image synthesis --- are wired in a graceful-degradation chain: if the external APIs are unavailable, the platform falls back to a local Pillow renderer and a rule-based text simplifier, ensuring functionality at zero API cost.

Functional testing across twelve test cases confirms end-to-end correctness for all implemented features. Image-based lesson generation completes in 3--18 seconds depending on the API path taken; ASL video generation completes in 20--45 seconds for short sentences. The modular architecture supports future integration of sign-recognition models, additional sign languages, and adaptive learning features. This work contributes to assistive educational technology and promotes inclusive learning for the hearing-impaired community.

\textbf{Keywords:} American Sign Language, Hearing-Impaired Education, WLASL Dataset, FastAPI, React, Google Gemini, Stability AI, Inclusive Technology.

\clearpage

% =====================================================================
% TABLE OF CONTENTS
% =====================================================================
\renewcommand{\contentsname}{TABLE OF CONTENTS}
\tableofcontents
\clearpage

\renewcommand{\listfigurename}{LIST OF FIGURES}
\listoffigures
\addcontentsline{toc}{chapter}{LIST OF FIGURES}
\clearpage

\renewcommand{\listtablename}{LIST OF TABLES}
\listoftables
\addcontentsline{toc}{chapter}{LIST OF TABLES}
\clearpage

% =====================================================================
% ABBREVIATIONS
% =====================================================================
\chapter*{ABBREVIATIONS}
\addcontentsline{toc}{chapter}{ABBREVIATIONS}

\begin{tabbing}
\hspace{3.5cm}\=\kill
\textbf{API}   \> Application Programming Interface \\
\textbf{ASL}   \> American Sign Language \\
\textbf{BSL}   \> British Sign Language \\
\textbf{CNN}   \> Convolutional Neural Network \\
\textbf{CORS}  \> Cross-Origin Resource Sharing \\
\textbf{CPU}   \> Central Processing Unit \\
\textbf{CRUD}  \> Create, Read, Update, Delete \\
\textbf{FFmpeg}\> Fast Forward Moving Picture Experts Group (multimedia framework) \\
\textbf{GAN}   \> Generative Adversarial Network \\
\textbf{GPU}   \> Graphics Processing Unit \\
\textbf{HCI}   \> Human--Computer Interaction \\
\textbf{ISL}   \> Indian Sign Language \\
\textbf{JSON}  \> JavaScript Object Notation \\
\textbf{JWT}   \> JSON Web Token \\
\textbf{LLM}   \> Large Language Model \\
\textbf{LMS}   \> Learning Management System \\
\textbf{LTI}   \> Learning Tools Interoperability \\
\textbf{MoSCoW}\> Must-Have, Should-Have, Could-Have, Won't-Have (prioritisation) \\
\textbf{NLP}   \> Natural Language Processing \\
\textbf{PNG}   \> Portable Network Graphics \\
\textbf{REST}  \> Representational State Transfer \\
\textbf{SDG}   \> Sustainable Development Goal \\
\textbf{SDXL}  \> Stable Diffusion Extra Large \\
\textbf{SPA}   \> Single-Page Application \\
\textbf{SQL}   \> Structured Query Language \\
\textbf{UI}    \> User Interface \\
\textbf{URL}   \> Uniform Resource Locator \\
\textbf{WCAG}  \> Web Content Accessibility Guidelines \\
\textbf{WLASL} \> Word-Level American Sign Language (dataset) \\
\end{tabbing}
\clearpage

% =====================================================================
% MAIN MATTER
% =====================================================================
\pagenumbering{arabic}
\setcounter{page}{1}

% =====================================================================
% CHAPTER 1 -- INTRODUCTION
% =====================================================================
\chapter{INTRODUCTION}

\section{Introduction to the Project}

The rapid advancement of digital technologies has significantly transformed modern education systems. However, accessibility remains a major challenge, especially for hearing-impaired learners who rely primarily on visual communication methods such as sign language. Most existing educational content is designed for text-based or audio-based consumption, which limits accessibility for deaf and hard-of-hearing children and creates a structural gap in inclusive education.

Storytelling and structured lesson delivery play a vital role in early childhood education by enhancing language development, imagination, creativity, and cognitive skills. Traditional methods do not adequately support hearing-impaired learners because they lack sign-language interpretation or intelligent simplification of complex sentence structures. As a result, these learners frequently struggle to engage with standard curricula.

This project addresses the gap directly through two intelligent pipelines. The first retrieves and assembles authentic ASL video clips from the WLASL dataset \cite{li2020wlasl} to produce continuous, natural sign-language storytelling videos. The second leverages Google Gemini AI \cite{team2023gemini} to simplify lesson text into concise, sign-friendly sequences and then triggers contextual image generation via Stability AI SDXL 1.0 \cite{podell2023sdxl} or a local Pillow-based renderer, so that educators can deliver structured, AI-enhanced lessons without requiring external hardware or human interpreters.

The full-stack system is accessible through any web browser and provides distinct user flows for educators (creating AI lessons) and learners (viewing ASL video and image content). This project aims to bridge the communication gap between textual educational content and visual learning, thereby promoting inclusive education for hearing-impaired children.

\subsection*{Platform Architecture Overview}
\begin{itemize}[itemsep=0.1em,topsep=0.2em,leftmargin=2em]
  \item \textbf{Text ingestion.} An educator enters lesson text and a title into the React frontend.
  \item \textbf{AI simplification.} The FastAPI backend calls Gemini 2.5 Flash to simplify the text into structured ASL-friendly sentences and word sequences, returned as JSON.
  \item \textbf{Image generation.} Each concept word drives a Stability AI SDXL prompt; if the API is unavailable, a local Pillow renderer generates colour-coded concept cards on-device.
  \item \textbf{ASL video assembly.} In parallel, the WLASL pipeline tokenises the text, retrieves matching clip files, standardises them via FFmpeg, and concatenates them with MoviePy into a single output video.
  \item \textbf{Result delivery.} The frontend renders the AI lesson structure, concept image cards, and ASL video on a single result page.
\end{itemize}

\begin{figure}[h]
\centering
\includegraphics[width=0.95\textwidth]{fig_home_page.png}
\caption{Platform home page showing the primary call-to-action for lesson creation.}
\label{fig:homepage}
\end{figure}

\section{Problem Statement}

Most digital educational tools treat accessibility for the hearing-impaired as an afterthought. Sign-language support is either absent entirely or provided as static vocabulary dictionaries with no continuous-narrative capability. When educational institutions attempt to provide ASL support, they depend on live human interpreters whose availability is limited, whose coverage is restricted to scheduled sessions, and whose cost is prohibitive for under-resourced schools.

Three concrete problems follow from this situation. First, hearing-impaired children cannot independently access the same lesson content their hearing peers receive, creating learning gaps that compound over time. Second, the available text-to-sign tools are either avatar-based systems that lack the natural expressiveness of real signers, or they require hardware (depth cameras, motion-capture suits) that most schools do not own. Third, no widely available platform combines real ASL video retrieval with Generative AI simplification and image synthesis; the two halves of the problem --- vocabulary mapping and concept visualisation --- are addressed separately if at all.

This project addresses the integration gap directly. The objective is to design and deploy a web application that converts lesson text into both a real-sign ASL video and a set of AI-generated concept images through one authentication layer, one input form, and one result page, so that an educator or learner performs the full workflow without switching between tools.

\section{Motivation}

The motivation for this platform arises from several converging observations.

\textbf{Scale of the gap.} Approximately 466 million people worldwide live with disabling hearing loss, of whom roughly 34 million are children \cite{who2021deafness}. For the vast majority, access to sign-language-interpreted educational content is limited or non-existent.

\textbf{Maturity of open datasets.} The WLASL dataset \cite{li2020wlasl}, released in 2020, provides over 21,000 video instances covering 2,000 ASL words performed by human signers. The How2Sign dataset \cite{duarte2021how2sign} adds continuous ASL narration over instructional videos. These resources make a data-driven, clip-retrieval approach practical without requiring a motion-capture recording studio.

\textbf{Advances in Generative AI.} Large language models such as Gemini \cite{team2023gemini} can reliably simplify complex sentences into short, sign-friendly word lists, reducing the vocabulary-mapping burden. Diffusion models such as Stability AI SDXL \cite{podell2023sdxl} can generate contextual educational images on demand, complementing the sign video with visual concept reinforcement --- a pedagogical technique well-supported in accessible education literature \cite{mayer2009multimedia}.

\textbf{Deployability on commodity hardware.} The architecture is designed to run on a single developer laptop or a free-tier cloud instance. The Pillow fallback ensures that educators in bandwidth-constrained environments can still generate concept cards without any external API call.

\section{Sustainable Development Goal of the Project}

The project aligns with three United Nations Sustainable Development Goals.

\textbf{SDG 4 -- Quality Education.} The platform ensures inclusive and equitable quality education by enabling hearing-impaired learners to access structured lesson content in their primary visual communication mode. It reduces the dependency on human interpreters and provides a scalable, always-available alternative.

\textbf{SDG 10 -- Reduced Inequalities.} By removing communication barriers and providing a freely deployable tool, the system minimises educational disparities between hearing-impaired learners and their hearing peers, particularly in under-resourced settings.

\textbf{SDG 9 -- Industry, Innovation, and Infrastructure.} The platform leverages cutting-edge AI (Gemini, Stability AI SDXL) and modern web technologies (FastAPI, React 18) to build assistive infrastructure that can be integrated into schools and e-learning platforms at minimal cost.

\section{Product Vision Statement}

Our vision is to enable inclusive and accessible learning experiences for hearing-impaired children by creating an intelligent, AI-powered ASL platform that transforms textual lesson content into meaningful sign-language narratives and contextual visual concept cards. We aim to bridge the communication gap between traditional educational resources and visual learners by leveraging real WLASL video data and Generative AI to deliver authentic, expressive, and engaging educational content.

We envision a digital ecosystem where educators, parents, and learners can seamlessly create and access lesson content in sign language --- ensuring that every child, regardless of hearing ability, can learn, imagine, and grow with confidence. Through innovation and inclusivity, this project strives to contribute toward equitable education and support global efforts in achieving the United Nations SDGs.

\section{Product Goal}

To develop and deploy an accessible, scalable, and user-friendly AI-powered learning platform that converts textual lesson content into ASL video narratives and AI-generated concept images, provides a responsive web interface for educators and learners, operates with graceful degradation when AI APIs are unavailable, and supports future extension with additional sign languages and recognition features.

% =====================================================================
% CHAPTER 2 -- LITERATURE SURVEY
% =====================================================================
\chapter{LITERATURE SURVEY}

\section{Overview of the Research Area}

The problem addressed in this project spans four intersecting research communities: sign language recognition and generation, natural language processing for accessible communication, Generative AI (large language models and diffusion models), and accessible educational technology. Sign language research appears primarily in computer vision venues (CVPR, ICCV, ECCV) and the ACM SIGACCESS conference on assistive technology. NLP for accessibility is spread across both mainstream NLP venues (ACL, EMNLP) and HCI conferences (CHI, ASSETS). Generative AI research, the fastest-moving of the four, is concentrated in arXiv preprints and NeurIPS/ICLR venues.

Despite the topical separation, the underlying methodological challenges are similar. All four fields grapple with data scarcity: annotated ASL corpora are expensive to collect; accessibility-focused NLP datasets are limited in size; and document-level sign-language datasets almost always under-represent the continuous, narratively coherent signing that educational storytelling requires. A second shared challenge is naturalness: avatar-based signing and rule-based simplification both produce output that trained signers and deaf learners find stilted. A third challenge is deployability: most published work reports offline metrics without attention to inference latency, memory footprint, or the experience of a non-technical user.

This chapter surveys the prior work that informed the platform's design, identifies the gaps that remain, states the research objectives that follow, and presents the product backlog and project roadmap.

\section{Existing Models and Frameworks}

\subsection{Sign Language Recognition}

Early sign language recognition relied on hand-crafted visual features --- hand-shape descriptors, wrist trajectories, and spatial position relative to the body --- combined with Hidden Markov Models for temporal modelling \cite{starner1995asl}. These approaches achieved reasonable performance on constrained vocabulary sets but failed to generalise to continuous signing with coarticulation and natural variation.

The deep learning era shifted the field toward CNNs for spatial feature extraction and RNNs or LSTMs for temporal modelling. Koller et~al.\ \cite{koller2015continuous} demonstrated that a CNN-LSTM hybrid could decode continuous German Sign Language at word-level accuracy approaching human transcription on a single-signer corpus. The I3D and SlowFast architectures, originally developed for action recognition, were subsequently adapted for sign recognition with competitive results on the WLASL benchmark \cite{li2020wlasl}.

Transformer-based models have become dominant since 2021. Camg\"oz et~al.\ \cite{camgoz2020sign} extended the neural machine translation paradigm to sign language translation, using a Transformer encoder over sign-sequence features to produce spoken-language text. More recent work, including the SignBERT family, applies self-supervised pre-training to hand-pose sequences before fine-tuning on labelled sign corpora.

\subsection{ASL Datasets}

The availability of large-scale ASL corpora has accelerated progress. The WLASL dataset \cite{li2020wlasl}, with 21,083 video instances over 2,000 signs performed by 119 signers, is the largest word-level ASL benchmark available. The How2Sign dataset \cite{duarte2021how2sign} provides 80 hours of continuous ASL signing over instructional cooking videos, paired with English transcriptions, making it suitable for text-to-sign generation research. The ASL Citizen dataset \cite{desai2023aslcitizen} is a recent community-sourced collection of 84,000 clips from over 50 signers, covering 2,731 words with diversity of age, dialect, and signing style.

Our platform uses WLASL as its primary clip library. Each word in the platform's vocabulary corresponds to one or more WLASL clip files; the pipeline selects the first available clip for a matched token, meaning that the system's effective vocabulary is bounded by the local WLASL subset that has been downloaded.

\subsection{Text-to-Sign Generation}

Converting natural language text into sign-language video is harder than the reverse problem. Karpouzis et~al.\ \cite{karpouzis2007avatar} built one of the earliest avatar-based Greek Sign Language generators, using a motion-capture corpus and rule-based synthesis. The gap between avatar output and natural signing remains a recognised limitation: deaf users consistently rate real-signer clips as more comprehensible and more natural than avatar-synthesised output \cite{bragg2019sign}.

Clip-retrieval approaches like ours sidestep the avatar naturalness problem by selecting real video clips for each sign and concatenating them. The main limitation is vocabulary coverage: any word not in the clip library must be fingerspelled, skipped, or approximated by a synonym. Prior systems using this approach include SignSynth \cite{efthimiou2004signsynth} and the platform described by Piao et~al.\ \cite{piao2020sign}, which used the BSL SignBank corpus for British Sign Language generation.

Diffusion-based sign synthesis is an emerging direction. Recent work has used pose-conditioned video diffusion to generate continuous signing from gloss sequences, but these models require significant compute at inference time and are not yet practical for deployment on free-tier hardware.

\subsection{AI in Accessible Education}

The use of AI to make educational content more accessible has grown substantially. Mayer's multimedia learning theory \cite{mayer2009multimedia} provides the pedagogical foundation: dual-channel processing (visual + verbal) improves comprehension for all learners, and is especially important for hearing-impaired learners for whom the verbal channel is inaccessible in its audio form. Replacing the audio channel with visual sign-language video and supplementary concept images directly maps onto Mayer's framework.

Large language models have been applied to educational content simplification with promising results. Shardlow \cite{shardlow2014text} surveyed classical text simplification approaches (lexical substitution, syntactic restructuring, content removal); modern LLMs can perform all three in a single generation step. Gemini 2.5 Flash is particularly well-suited for this task because its instruction-following capability is strong enough to produce consistently structured JSON output, making downstream parsing reliable.

\subsection{Generative AI for Visual Content}

Diffusion models have transformed AI image generation since the release of Stable Diffusion in 2022 \cite{rombach2022high}. The SDXL 1.0 model \cite{podell2023sdxl} produces 1024$\times$1024 images of sufficient quality for educational concept illustration from short text prompts. For this platform, prompts are structured to emphasise educational clarity: white background, educational style, and explicit naming of the sign concept. When the API is unavailable, the platform's Pillow-based fallback generates 512$\times$512 concept cards locally using Python's Pillow library, with colour-coded rounded borders, an ASL badge, and the concept word in large type.

\section{Limitations Identified from Literature Survey (Research Gaps)}

The survey above surfaces several gaps that informed the design of this platform.

\textbf{G1: Naturalness vs.\ coverage trade-off is unresolved.} Real-clip retrieval produces natural signs but is limited to the vocabulary in the local dataset. Avatar synthesis covers arbitrary vocabulary but produces output deaf learners find less comprehensible. No deployed system resolves this trade-off for general educational use.

\textbf{G2: AI-aided lesson simplification is underexplored for ASL.} Most ASL tools map text to signs at the word level without preprocessing the source text. A sentence like ``Water evaporates from the ocean surface'' cannot be signed word-for-word using a clip-retrieval system; intelligent simplification to ``water evaporate ocean'' is required but rarely automated.

\textbf{G3: Visual concept reinforcement is absent from existing ASL platforms.} Published ASL education tools focus exclusively on sign video and do not pair each concept with a contextual image. Mayer's dual-channel theory suggests this omission reduces comprehension, especially for abstract concepts.

\textbf{G4: Deployability is under-reported.} Most text-to-sign research reports only offline correctness metrics. Latency, memory footprint, and the ability to serve predictions inside a synchronous web request without a GPU are rarely addressed.

\textbf{G5: Graceful degradation is absent.} If an ASL platform's video generation fails, most systems return an error. A fallback chain that degrades gracefully --- from AI-generated images to locally rendered concept cards to lesson structure only --- is not described in the literature.

\section{Research Objectives}

The gaps above translate into the following research objectives.

\begin{enumerate}[leftmargin=2em,itemsep=0.2em]
  \item Design and implement a clip-retrieval ASL video pipeline using the WLASL dataset, evaluated on a sample of WLASL vocabulary and confirmed to produce smooth, natural-looking sign-language video.
  \item Integrate Google Gemini 2.5 Flash as an AI lesson simplification layer, evaluating its ability to convert arbitrary lesson text into structured ASL-friendly JSON with original, simplified, and ASL-sequence fields.
  \item Implement a Stability AI SDXL image generation pipeline for concept visualisation, and a local Pillow fallback that produces concept cards with zero API dependency.
  \item Deploy the full stack as a FastAPI + React 18 web application within a strict budget: no GPU, no paid hosting, and functional latency on a developer laptop.
  \item Document the performance and limitations of each pipeline component honestly, so subsequent work has a reproducible baseline.
\end{enumerate}

\section{Product Backlog (User Stories with Desired Outcomes)}

The objectives above were broken down into user stories during sprint planning. Table~\ref{tab:backlog} summarises the product backlog at the start of the project.

\begin{table}[h]
\caption{Product backlog at project kick-off}
\label{tab:backlog}
\centering
\small
\begin{tabularx}{\linewidth}{@{}lXX@{}}
\toprule
\textbf{ID} & \textbf{User Story} & \textbf{Desired Outcome} \\
\midrule
US-01 & As a learner, I want to watch a continuous ASL video for a given lesson so that I can understand the content in sign language. & Concatenated sign video returned within 60 seconds for a short sentence; smooth playback in the browser. \\
US-02 & As an educator, I want AI to simplify my lesson text so that the signs are accurate and the sequence is natural. & Gemini returns a structured JSON with simplified sentences and ASL word sequences. \\
US-03 & As a learner, I want to see visual concept images for each lesson word so that my comprehension is reinforced beyond the sign video. & AI-generated or locally rendered concept card for each word, displayed in a responsive grid. \\
US-04 & As a user, I want to register and log in so that my lessons are saved and accessible on future visits. & JWT-based authentication working end-to-end; lessons persisted to SQLite. \\
US-05 & As a user, I want a clean and responsive result page so that the lesson content is easy to navigate. & React result page loads without errors; all hooks comply with React Rules of Hooks. \\
US-06 & As a developer, I want unknown ASL words handled gracefully so that the system does not fail on missing clips. & Unknown words skipped; pipeline continues with available clips; no 500 error. \\
US-07 & As a system, I want the image pipeline to degrade gracefully so that lesson images are always generated. & If Stability AI fails, Pillow cards are generated locally; lesson returned regardless. \\
US-08 & As an operator, I want all file paths to be dynamic so that the system works on any machine. & No hardcoded paths in any service file; all paths derived from \texttt{\_\_file\_\_} or environment variables. \\
\bottomrule
\end{tabularx}
\end{table}

\section{Plan of Action (Project Roadmap)}

The project ran over two formal sprints between December 2025 and April 2026. Table~\ref{tab:roadmap} lists the planned milestones with the actual delivery periods.

\begin{table}[h]
\caption{Project roadmap and milestones}
\label{tab:roadmap}
\centering
\small
\begin{tabularx}{\linewidth}{@{}lXl@{}}
\toprule
\textbf{Phase} & \textbf{Milestone} & \textbf{Period} \\
\midrule
Setup     & Repo scaffolding, FastAPI skeleton, React+Vite shell, SQLite schema, WLASL clip setup & Wk 1--4 \\
Sprint I  & WLASL video pipeline, JWT authentication, basic React UI (Home, SignIn, CreateLesson, VideoResult) & Wk 5--12 \\
Mid-term review & Faculty panel demo; backlog re-prioritised based on feedback & Wk 13 \\
Sprint II & Gemini AI lesson service, Stability AI image generation, Pillow fallback, AI lesson result page & Wk 14--22 \\
Hardening & Path normalisation across all services, bug fixes, CORS and hooks issues resolved & Wk 23--24 \\
Documentation & Project report, slide deck, plagiarism check & Wk 25--26 \\
\bottomrule
\end{tabularx}
\end{table}

% =====================================================================
% CHAPTER 3 -- SPRINT PLANNING AND EXECUTION METHODOLOGY
% =====================================================================
\chapter{SPRINT PLANNING AND EXECUTION METHODOLOGY}

The project followed a two-sprint agile cadence with a faculty review at the boundary. Sprint I produced an authenticated single-page application with a working ASL video pipeline and a minimal result page. Sprint II added the Gemini AI lesson service, the Stability AI and Pillow image pipelines, and a polished result page that displays both images and video. This chapter documents the objectives, functional and architectural artefacts, results, and retrospective notes for each sprint.

\section{Sprint I}

\subsection{Objectives with User Stories of Sprint I}

Sprint I covered weeks 5 through 12. The objective was to ship a deployable application with a working end-to-end ASL video pipeline and the supporting infrastructure (JWT auth, React routing, result page) that the Sprint II AI features could plug into. The user stories pulled into Sprint I were:

\begin{itemize}[leftmargin=2em,itemsep=0.1em]
  \item \textbf{US-04} -- JWT-based authentication, register and login.
  \item \textbf{US-01} -- WLASL clip retrieval and video concatenation pipeline.
  \item \textbf{US-06} -- Graceful handling of unknown ASL words.
  \item \textbf{US-08} -- Dynamic file paths; no hardcoded machine-specific directories.
\end{itemize}

The acceptance criterion for the sprint was a locally running FastAPI + React application that accepted lesson text, generated an ASL video, and played it back in the browser --- end to end --- without requiring manual path edits.

\subsection{Functional Document}

The functional scope of Sprint I covered the following capabilities, each tied to a user story.

\textbf{F-01: User authentication.} Email-and-password sign-up and sign-in via a FastAPI auth router. Passwords are hashed with bcrypt and stored in SQLite via SQLAlchemy. The \texttt{/auth/register} and \texttt{/auth/login} endpoints return signed JWT tokens that the React frontend stores in \texttt{localStorage} and attaches to all subsequent requests.

\textbf{F-02: Text-to-ASL video pipeline.} The \texttt{dynamic\_asl\_generator.py} module accepts a text string via the \texttt{/generate-asl} endpoint. It tokenises and lowercases the input, looks up each token against a JSON vocabulary index derived from WLASL, retrieves the matching \texttt{.mp4} clip from \texttt{storage/asl\_clips/}, standardises clip resolution and frame rate via FFmpeg, and concatenates clips with MoviePy 1.0.3 into a final video saved to \texttt{storage/processed/}. Unknown tokens are silently skipped; the pipeline continues with available clips.

\textbf{F-03: Static file serving.} FastAPI mounts the \texttt{storage/} directory at the \texttt{/static} path so that generated videos are accessible to the React frontend by URL without a separate file server.

\textbf{F-04: Basic React UI.} A React 18 + Vite single-page application provides a Home page, a Sign In / Sign Up page with form validation, a Create Lesson page with a text area and title field, and a Video Result page that streams the generated ASL video.

\subsection{Architecture Document}

The Sprint I architecture established the three-tier layout that the rest of the project built upon. The presentation tier was a React 18 SPA built with Vite 5. The application tier was a FastAPI service exposing RESTful endpoints with Pydantic request/response validation and a JWT authentication middleware. The storage tier was a SQLite database (via SQLAlchemy) for user and lesson metadata, plus a local file system for ASL clips and generated videos.

All file paths were made dynamic from the start. Every service module derived its root directory from \texttt{os.path.abspath(os.path.join(os.path.dirname(\_\_file\_\_), \ldots))}, preventing the hardcoded-path failures that plagued early development.

\begin{figure}[h]
\centering
\includegraphics[width=0.95\textwidth]{fig_architecture.png}
\caption{System architecture: React 18 SPA communicating with the FastAPI gateway, which orchestrates the WLASL video pipeline and AI services.}
\label{fig:architecture}
\end{figure}

\begin{figure}[h]
\centering
\includegraphics[width=0.48\textwidth]{fig_signin_page.png}\hfill
\includegraphics[width=0.48\textwidth]{fig_signup_page.png}
\caption{Authentication pages: Sign-In (left) and Sign-Up (right).}
\label{fig:auth}
\end{figure}

\begin{figure}[h]
\centering
\includegraphics[width=0.48\textwidth]{fig_create_lesson.png}\hfill
\includegraphics[width=0.48\textwidth]{fig_asl_pipeline.png}
\caption{Create Lesson page (left) and the WLASL clip-retrieval and concatenation pipeline diagram (right).}
\label{fig:sprint1_ui}
\end{figure}

\textbf{Technology Stack -- Sprint I:}

\begin{center}
\small
\begin{tabular}{ll}
\toprule
\textbf{Component} & \textbf{Technology} \\
\midrule
Backend Framework & FastAPI (Python 3.11) \\
Frontend Framework & React 18 + Vite 5 \\
Database & SQLite via SQLAlchemy \\
Authentication & JWT (python-jose, bcrypt) \\
Video Processing & MoviePy 1.0.3 + FFmpeg \\
ASL Dataset & WLASL (Word-Level ASL) \\
HTTP Client & Axios (React frontend) \\
\bottomrule
\end{tabular}
\end{center}

\subsection{Outcome of Objectives / Result Analysis}

By the end of Sprint I, all four committed user stories were delivered. The WLASL video pipeline correctly retrieved and concatenated clips for all words present in the vocabulary index. The JWT authentication flow worked end-to-end; tokens were validated on protected routes and rejected on expiry. The React result page streamed the generated video without buffering issues on a local network.

Two issues surfaced during Sprint I and were fixed before the mid-term review. First, MoviePy version 2 had removed the \texttt{moviepy.editor} import path used in the original generator; the fix was to pin \texttt{moviepy==1.0.3} in \texttt{requirements.txt}. Second, several service files contained hardcoded absolute paths to a specific developer's machine (\texttt{/Users/arihantjain/Desktop/\ldots}); these were replaced with dynamic path calculations using \texttt{\_\_file\_\_}-relative expressions.

\subsection{Sprint Retrospective}

\textit{What went well.} The decision to treat authentication and routing as Sprint I deliverables paid off; integrating the Sprint II AI features required no frontend structural changes. The dynamic path strategy eliminated environment-specific setup entirely.

\textit{What did not go well.} MoviePy's major-version API change (v1 vs.\ v2) was not anticipated and cost several hours of debugging. The WLASL vocabulary index required manual curation to remove clips whose file paths did not match the downloaded subset.

\textit{What we learned.} Pin dependency versions explicitly from the start. Verify that every vocabulary entry in a dataset index actually corresponds to a present file before wiring the index into the pipeline.

\textit{What to change next sprint.} Wire the AI lesson service as a new FastAPI router (not a modification of the video router) to maintain clean module separation. Add a \texttt{/health} endpoint that reports which services are available.

\section{Sprint II}

\subsection{Objectives with User Stories of Sprint II}

Sprint II covered weeks 14 through 22. The objective was to deliver the Gemini AI lesson simplification service, the Stability AI and Pillow image generation pipelines, and a polished result page that displays both images and structured lesson data. The user stories pulled into Sprint II were:

\begin{itemize}[leftmargin=2em,itemsep=0.1em]
  \item \textbf{US-02} -- Gemini AI lesson simplification.
  \item \textbf{US-03} -- Stability AI image generation and Pillow fallback.
  \item \textbf{US-05} -- Polished AI lesson result page (React hooks compliance, image grid).
  \item \textbf{US-07} -- Graceful degradation chain for image generation.
\end{itemize}

\subsection{Functional Document}

\textbf{F-05: Gemini AI lesson simplification (\texttt{ai\_lesson\_service.py}).} The \texttt{AILessonService} class calls \texttt{gemini-2.5-flash} with a structured prompt requesting JSON output. Each sentence is decomposed into \texttt{original}, \texttt{simplified}, and \texttt{asl\_sequence} fields. A regex-based JSON extractor parses the model response. If Gemini is unavailable (no API key or quota exceeded), the service falls back to a rule-based simplifier that removes stopwords and limits each sequence to four key words.

\textbf{F-06: Stability AI SDXL image generation (\texttt{working\_image\_service.py}).} The \texttt{WorkingImageService} calls the Stability AI REST API at \texttt{stable-diffusion-xl-1024-v1-0/text-to-image} with a prompt structured for educational clarity. The response contains a base64-encoded PNG that is decoded and saved to \texttt{storage/processed/images/}. SSL certificate verification is disabled (\texttt{verify=False}) to work around a certificate chain issue on macOS with Python 3.14; urllib3 warnings are suppressed accordingly.

\textbf{F-07: Pillow fallback concept cards (\texttt{fallback\_image\_service.py}).} When Stability AI is unavailable, the platform generates local ASL concept cards using Python's Pillow library. Each card is a 512$\times$512 PNG with a colour-coded rounded border, an ``ASL SIGN'' badge, a hand-icon placeholder, and the concept word in large type. A palette of six colours cycles across cards for visual variety. Cards are saved to the same \texttt{storage/processed/images/} path, making the frontend agnostic to the image source.

\begin{figure}[h]
\centering
\includegraphics[width=0.55\textwidth]{fig_asl_card.png}
\caption{Sample Pillow-generated ASL concept card (local fallback renderer).}
\label{fig:aslcard}
\end{figure}

\textbf{F-08: AI lesson FastAPI route (\texttt{routes/ai\_lesson.py}).} The \texttt{/ai-lesson/generate-asl-lesson} endpoint calls \texttt{ai\_lesson\_service.generate\_ai\_image\_lesson()}, which orchestrates the Gemini simplification, the image generation chain, and the WLASL lookup. It returns a structured \texttt{LessonResponse} Pydantic model containing the lesson data and image metadata.

\textbf{F-09: Polished result page (\texttt{VideoResultPage.jsx}).} The React result page was refactored to display the AI lesson structure (original and simplified sentences, ASL word sequences) alongside the concept image grid and the ASL video player. All React hooks (\texttt{useState}, \texttt{useMemo}, \texttt{useEffect}, \texttt{useRef}) were moved before any conditional early returns to comply with React's Rules of Hooks \cite{react2023hooks}, resolving a ``rendered more hooks than during the previous render'' runtime error that caused blank pages on lesson load.

\subsection{Architecture Document}

Sprint II extended the Sprint I architecture with three new components. The Gemini API and the Stability AI API were added as external dependencies, each wrapped in a stateless service class. The Pillow renderer was added as a local fallback module with no network dependency.

\begin{figure}[h]
\centering
\includegraphics[width=0.95\textwidth]{fig_architecture_sprint2.png}
\caption{Updated Sprint II architecture showing the Gemini and Stability AI external APIs alongside the Pillow local fallback.}
\label{fig:architecture2}
\end{figure}

The image generation pipeline implements a graceful-degradation chain. First, the platform attempts Stability AI SDXL. If the API key is absent or the call fails, it falls back to the Pillow local renderer. If the Pillow render also fails, the lesson structure is returned without images. The frontend handles all three outcomes without showing an error to the user.

\begin{figure}[h]
\centering
\includegraphics[width=0.95\textwidth]{fig_result_page.png}
\caption{AI lesson result page showing the simplified lesson structure, concept image grid, and ASL video player.}
\label{fig:resultpage}
\end{figure}

\textbf{Additional Technology added in Sprint II:}

\begin{center}
\small
\begin{tabular}{ll}
\toprule
\textbf{Component} & \textbf{Technology} \\
\midrule
AI Text Simplification & Google Gemini 2.5 Flash \\
AI Image Generation & Stability AI SDXL 1.0 \\
Fallback Image Renderer & Python Pillow 10.x \\
HTTP requests (backend) & Python requests library \\
\bottomrule
\end{tabular}
\end{center}

\subsection{Outcome of Objectives / Result Analysis}

All four Sprint II user stories were delivered. The Gemini service correctly returned structured JSON for all test lesson inputs, including edge cases with long sentences, numbers, and domain-specific vocabulary. The Stability AI pipeline returned 1024$\times$1024 PNG concept images for each keyword. When the API key was removed from \texttt{.env}, the Pillow fallback engaged automatically and produced concept cards within one second. The React result page rendered without hooks errors after the refactor.

Two issues surfaced during Sprint II. First, an SSL certificate chain error on macOS with Python 3.14 prevented Stability AI requests from completing; this was resolved with \texttt{verify=False} and urllib3 warning suppression. Second, a CORS middleware crash caused by whitespace in the \texttt{CORS\_ORIGINS} environment variable was fixed by stripping each origin string at startup.

\subsection{Sprint Retrospective}

\textit{What went well.} The Gemini integration was straightforward; the model's structured JSON output was reliable enough that the regex extractor rarely triggered the fallback. The Pillow fallback means the platform has zero external API dependency for basic usage.

\textit{What did not go well.} The React hooks ordering violation required two debugging rounds before the root cause was identified; the second round was needed because the \texttt{useEffect} for video event listeners was declared after the conditional returns, having been missed in the first pass. SSL certificate handling on macOS Python 3.14 was opaque; standard remedies (\texttt{certifi}, \texttt{SSL\_CERT\_FILE}) did not resolve the issue.

\textit{What we learned.} React's Rules of Hooks must be enforced at code-review time, not discovered at runtime. A linter rule (\texttt{eslint-plugin-react-hooks}) would have caught both violations immediately.

\textit{What to change next.} Add the ESLint hooks plugin to the frontend toolchain. Investigate a proper SSL certificate bundle for the Stability AI call rather than relying on \texttt{verify=False} in production.

% =====================================================================
% CHAPTER 4 -- RESULTS AND DISCUSSIONS
% =====================================================================
\chapter{RESULTS AND DISCUSSIONS}

\section{Project Outcomes (Performance Evaluation, Testing Results)}

This chapter consolidates the evaluation of the ASL AI Education Platform across all pipeline components. For each component, we report correctness on the functional test suite, end-to-end latency on a developer laptop, and a brief discussion of failure modes and their drivers.

\subsection{ASL Video Generation}

The WLASL clip-retrieval pipeline correctly tokenised and normalised input lesson text, matched tokens against the vocabulary index, retrieved the corresponding \texttt{.mp4} clips, standardised each clip to 720p/25fps via FFmpeg, and concatenated them into a single output video using MoviePy 1.0.3. Videos were served via FastAPI's \texttt{StaticFiles} mount and streamed directly in the React \texttt{<video>} element.

Table~\ref{tab:video_test} summarises the test cases for the video pipeline. All six test cases passed.

\begin{table}[h]
\caption{Functional test cases -- ASL video pipeline}
\label{tab:video_test}
\centering
\small
\begin{tabularx}{\linewidth}{@{}lXXl@{}}
\toprule
\textbf{TC} & \textbf{Input} & \textbf{Expected Output} & \textbf{Result} \\
\midrule
TC-01 & ``hello world'' & Video with two sign clips & Pass \\
TC-02 & ``water evaporate sky'' & Three-clip concatenated video & Pass \\
TC-03 & ``quantum entanglement'' & Empty or partial video; no crash & Pass \\
TC-04 & Empty string & 422 validation error & Pass \\
TC-05 & 10-word sentence (mixed known/unknown) & Video for known words only & Pass \\
TC-06 & Video URL returned & Video streams in \texttt{<video>} element & Pass \\
\bottomrule
\end{tabularx}
\end{table}

\subsection{Gemini AI Lesson Simplification}

The Gemini 2.5 Flash model successfully simplified lesson text into structured ASL-friendly JSON in all test cases. A representative example is shown below.

\textit{Input:} ``Water evaporates from the surface of the ocean and rises to form clouds in the sky, which then release rain.''

\textit{Output:}
\begin{lstlisting}[language={},frame=single,caption=Sample Gemini AI JSON response]
{
  "lesson_title": "The Water Cycle",
  "sentences": [
    {
      "original": "Water evaporates from the surface of the ocean",
      "simplified": "water evaporate ocean surface",
      "asl_sequence": ["water", "evaporate", "ocean", "surface"]
    },
    {
      "original": "and rises to form clouds in the sky",
      "simplified": "water form clouds sky",
      "asl_sequence": ["water", "form", "clouds", "sky"]
    },
    {
      "original": "which then release rain",
      "simplified": "clouds release rain",
      "asl_sequence": ["clouds", "release", "rain"]
    }
  ]
}
\end{lstlisting}

When the Gemini API key was absent, the rule-based fallback engaged and produced simplified sequences by removing stopwords and limiting each sequence to four key words. Fallback output was less semantically accurate but never caused a service error.

\subsection{Image Generation}

When a valid Stability AI API key was configured, the platform generated 1024$\times$1024 PNG concept images in 12--18 seconds per lesson (three images). When the API key was absent or the call failed, the Pillow fallback produced 512$\times$512 locally rendered concept cards in under two seconds. Both output paths returned image metadata in the same format, making the frontend agnostic to the source.

\begin{figure}[h]
\centering
\includegraphics[width=0.95\textwidth]{fig_image_pipeline.png}
\caption{Image generation pipeline showing the graceful-degradation chain: Stability AI SDXL $\rightarrow$ Pillow local renderer $\rightarrow$ lesson structure only.}
\label{fig:imagepipeline}
\end{figure}

\subsection{End-to-End Latency}

\begin{table}[h]
\caption{End-to-end latency on a developer laptop (Apple M1, 8 GB RAM)}
\label{tab:latency}
\centering
\small
\begin{tabular}{lc}
\toprule
\textbf{Operation} & \textbf{Latency} \\
\midrule
Initial page load (Home) & $<$\,1.5 s \\
ASL video generation (5-word sentence) & 20--45 s \\
AI lesson (Gemini + Stability AI, 3 images) & 12--18 s \\
AI lesson (Gemini + Pillow fallback, 3 images) & 3--5 s \\
AI lesson (rule-based + Pillow, no API keys) & $<$\,2 s \\
\bottomrule
\end{tabular}
\end{table}

The ASL video pipeline latency is dominated by FFmpeg clip normalisation and MoviePy concatenation, not by file I/O. The Stability AI latency is network-bound. The Pillow fallback is CPU-bound and completes well within the synchronous request window.

\subsection{Accessibility Review}

The platform was reviewed against WCAG 2.1 AA guidelines. All interactive elements carry ARIA labels. The colour contrast ratio on primary UI elements exceeds 4.5:1. The video player exposes keyboard controls via the native \texttt{<video>} element. The concept card grid uses \texttt{alt} attributes on all \texttt{<img>} elements with the concept word as the alt text.

\section{Committed vs Completed User Stories}

\begin{table}[h]
\caption{Committed vs completed user stories across both sprints}
\label{tab:committed}
\centering
\small
\begin{tabularx}{\linewidth}{@{}lXlll@{}}
\toprule
\textbf{ID} & \textbf{User Story} & \textbf{Committed} & \textbf{Completed} & \textbf{Sprint} \\
\midrule
US-01 & WLASL clip retrieval and video generation & Yes & Yes & Sprint I \\
US-04 & JWT authentication (register / login) & Yes & Yes & Sprint I \\
US-06 & Graceful unknown-word handling & Yes & Yes & Sprint I \\
US-08 & Dynamic file paths on all services & Yes & Yes & Sprint I \\
US-02 & Gemini AI lesson simplification & Yes & Yes & Sprint II \\
US-03 & Stability AI + Pillow image generation & Yes & Yes & Sprint II \\
US-05 & Polished result page (hooks compliance) & Yes & Yes & Sprint II \\
US-07 & Graceful degradation chain & Yes & Yes & Sprint II \\
Real-time sign recognition & Could have & No & -- \\
\bottomrule
\end{tabularx}
\end{table}

% =====================================================================
% CHAPTER 5 -- CONCLUSION AND FUTURE ENHANCEMENT
% =====================================================================
\chapter{CONCLUSION AND FUTURE ENHANCEMENT}

\section{Summary}

The ASL-Based Storytelling and AI-Powered Learning Platform is a working demonstration that a real-clip sign-language video pipeline and a Generative AI lesson generation pipeline can be integrated into one web application and delivered within a strict deployment budget. The platform successfully converts lesson text into both authentic ASL video and AI-generated concept images, with a graceful-degradation chain that produces useful output even with no paid API keys.

The project also provided concrete experience with the operational realities of deployed AI systems: SSL certificate handling on modern macOS, React Rules of Hooks enforcement, MoviePy major-version API breaks, hardcoded-path failures across developer environments, and the trade-offs between Stability AI image quality and local Pillow rendering speed. These lessons are reflected throughout the architecture and are documented in the sprint retrospectives above.

\section{Limitations}

Five limitations should be stated plainly.

\textit{Vocabulary coverage.} The WLASL clip-retrieval approach is bounded by the local clip subset. Words not in the vocabulary are silently skipped; a fingerspelling fallback or synonym substitution would improve coverage but adds complexity.

\textit{SSL certificate handling.} The \texttt{verify=False} workaround for Stability AI calls on macOS Python 3.14 disables certificate validation and should not be used in a public production deployment.

\textit{Gemini JSON reliability.} The regex-based JSON extractor occasionally falls back to the rule-based simplifier when the model produces well-formed but slightly offset JSON. A stricter function-calling interface (Gemini's structured output mode) would be more reliable.

\textit{Latency of video generation.} ASL video generation takes 20--45 seconds for a short sentence; this is dominated by FFmpeg and MoviePy processing. For a classroom deployment, pre-caching common words or switching to a streaming assembly approach would improve the user experience.

\textit{No adversarial or robustness testing.} The platform has not been tested under concurrent load, with malformed inputs beyond the validation layer, or against prompt-injection attacks on the Gemini endpoint.

\section{Future Enhancements}

Future work falls into three categories.

\textit{Modelling.} We plan to add a MediaPipe-based real-time ASL sign recognition module that allows learners to practise signs via webcam and receive instant feedback. Expanding the vocabulary index with ASL Citizen \cite{desai2023aslcitizen} clips for under-represented words is also planned, as is a fingerspelling fallback for out-of-vocabulary tokens.

\textit{Multilinguality.} Extending the platform beyond ASL to British Sign Language (BSL) and Indian Sign Language (ISL) would significantly broaden the user base. The architecture is modular enough to support a vocabulary-index swap per sign language.

\textit{Systems.} Replacing the \texttt{verify=False} Stability AI workaround with a properly configured certificate bundle is the immediate security priority. An asynchronous job queue for video generation would allow the frontend to poll for progress rather than blocking on a long synchronous request. An ESLint React hooks lint rule should be added to the CI pipeline.

\textit{Product.} An LTI (Learning Tools Interoperability) connector would allow the platform to be embedded in popular Learning Management Systems such as Moodle and Canvas. A React Native mobile application would make the platform accessible without a browser, which is important for learners in environments where laptops are scarce.

\section{Concluding Remarks}

The pace at which Generative AI capabilities are improving means that both the text simplification and the image generation components of this platform can be upgraded with minimal architectural change. The clip-retrieval pipeline is similarly replaceable: as sign-synthesis models mature and become deployable on commodity hardware, they can substitute for or augment the WLASL clips behind the same API contract. The bet underlying this platform is that integration --- one workflow, one history, one consistent way of delivering sign-language lesson content --- is more durable than any individual model version inside it.

% =====================================================================
% APPENDIX
% =====================================================================
\appendix

\chapter{SAMPLE CODE}

\section{FastAPI AI Lesson Endpoint}

\begin{lstlisting}[language=Python,caption={AI Lesson Route (app/routes/ai\_lesson.py)}]
@router.post("/generate-asl-lesson", response_model=LessonResponse)
async def generate_asl_lesson(request: LessonRequest,
                               background_tasks: BackgroundTasks):
    ai_available = is_ai_available()
    try:
        lesson_data = await ai_lesson_service.generate_ai_image_lesson(
            request.lesson_text, request.lesson_title)
        if lesson_data.get("ai_images_used") and lesson_data.get("image_data"):
            return LessonResponse(
                success=True,
                lesson_data=lesson_data,
                image_data=lesson_data["image_data"],
                message="ASL lesson with images generated successfully!",
                ai_available=ai_available)
        return LessonResponse(
            success=True, lesson_data=lesson_data,
            message="Lesson generated (image generation unavailable)",
            ai_available=ai_available)
    except Exception as e:
        raise HTTPException(status_code=500,
                            detail=f"Failed to generate ASL lesson: {e}")
\end{lstlisting}

\section{Gemini AI Lesson Service}

\begin{lstlisting}[language=Python,caption={Gemini simplification and JSON parsing (ai\_lesson\_service.py)}]
def generate_asl_lesson(self, lesson_text: str,
                         lesson_title: str = "Untitled") -> Dict:
    if not self.is_available():
        return self._fallback_generation(lesson_text, lesson_title)
    prompt = self._build_prompt(lesson_text, lesson_title)
    response = self.model.generate_content(prompt)
    return self._parse_ai_response(response.text, lesson_text, lesson_title)

def _parse_ai_response(self, response_text: str,
                        original_text: str, lesson_title: str) -> Dict:
    json_match = re.search(r'\{.*\}', response_text, re.DOTALL)
    if json_match:
        data = json.loads(json_match.group())
        if "lesson_title" in data and "sentences" in data:
            return data
    return self._fallback_generation(original_text, lesson_title)
\end{lstlisting}

\section{Pillow Fallback Concept Card}

\begin{lstlisting}[language=Python,caption={Local ASL concept card renderer (fallback\_image\_service.py)}]
def _draw_card(concept: str, image_number: int,
               size: int = 512) -> Image.Image:
    fg, bg = _PALETTE[(image_number - 1) % len(_PALETTE)]
    img = Image.new("RGB", (size, size), bg)
    draw = ImageDraw.Draw(img)
    draw.rounded_rectangle([12, 12, size-12, size-12],
                            radius=32, outline=fg, width=6)
    try:
        font = ImageFont.truetype(
            "/System/Library/Fonts/Helvetica.ttc", 48)
    except Exception:
        font = ImageFont.load_default()
    lw = draw.textlength(concept.upper(), font=font)
    draw.text(((size - lw) / 2, size - 110),
              concept.upper(), font=font, fill=fg)
    return img
\end{lstlisting}

\section{React Result Page -- Hooks Before Returns}

\begin{lstlisting}[language=JavaScript,caption={VideoResultPage.jsx -- all hooks before conditional returns}]
export default function VideoResultPage() {
  const [status, setStatus] = useState(null);
  const [loading, setLoading] = useState(true);
  const [error, setError]   = useState(null);
  const videoRef = useRef(null);

  // ALL hooks declared BEFORE any conditional return
  const teacherTokens = useMemo(() => computeTokens(status), [status]);
  const aslVideoSrc   = useMemo(() => buildVideoUrl(status), [status]);

  useEffect(() => {
    if (!videoRef.current) return;
    const vid = videoRef.current;
    const onErr = () => console.error("Video error:", aslVideoSrc);
    vid.addEventListener("error", onErr);
    return () => vid.removeEventListener("error", onErr);
  }, [aslVideoSrc]);

  // Conditional returns AFTER all hooks
  if (loading) return <LoadingSpinner />;
  if (error)   return <ErrorMessage msg={error} />;
  // ... render lesson content
}
\end{lstlisting}

% =====================================================================
% REFERENCES
% =====================================================================
\renewcommand\bibname{REFERENCES}
\addcontentsline{toc}{chapter}{REFERENCES}

\begin{thebibliography}{30}

\bibitem{li2020wlasl}
D.\ Li, C.\ Rodriguez, X.\ Yu, and H.\ Li,
``Word-level deep sign language recognition from video: A new large-scale dataset and methods comparison,''
\textit{Proc.\ IEEE/CVF WACV}, pp.\ 1459--1469, 2020.

\bibitem{team2023gemini}
R.\ Anil, A.\ Chowdhery, J.\ Dean \textit{et al.},
``Gemini: A family of highly capable multimodal models,''
\textit{arXiv preprint arXiv:2312.11805}, 2023.

\bibitem{podell2023sdxl}
D.\ Podell, Z.\ English, K.\ Lacey, A.\ Blattmann, T.\ Dockhorn, J.\ M\"uller, J.\ Penna, and R.\ Rombach,
``SDXL: Improving latent diffusion models for high-resolution image synthesis,''
\textit{arXiv preprint arXiv:2307.01952}, 2023.

\bibitem{rombach2022high}
R.\ Rombach, A.\ Blattmann, D.\ Lorenz, P.\ Esser, and B.\ Ommer,
``High-resolution image synthesis with latent diffusion models,''
\textit{Proc.\ IEEE/CVF CVPR}, pp.\ 10684--10695, 2022.

\bibitem{duarte2021how2sign}
A.\ Duarte, S.\ Palaskar, L.\ Ventura, D.\ Ghadiyaram, K.\ DeHaan, F.\ Metze, J.\ Torres, and X.\ Giro-i-Nieto,
``How2Sign: A large-scale multimodal dataset for continuous American Sign Language,''
\textit{Proc.\ IEEE/CVF CVPR}, pp.\ 2241--2250, 2021.

\bibitem{desai2023aslcitizen}
A.\ Desai, L.\ Berger, M.\ Bhatt \textit{et al.},
``ASL Citizen: A community-sourced dataset for advancing isolated sign language recognition,''
\textit{arXiv preprint arXiv:2304.05934}, 2023.

\bibitem{who2021deafness}
World Health Organization,
``Deafness and hearing loss,''
WHO Fact Sheet, 2021.

\bibitem{mayer2009multimedia}
R.\ E.\ Mayer,
\textit{Multimedia Learning}, 2nd ed.
Cambridge University Press, 2009.

\bibitem{bragg2019sign}
D.\ Bragg, O.\ Koller, M.\ Bellard \textit{et al.},
``Sign language recognition, generation, and translation: An interdisciplinary perspective,''
\textit{Proc.\ ACM SIGACCESS ASSETS}, pp.\ 16--31, 2019.

\bibitem{camgoz2020sign}
N.\ C.\ Camg\"oz, O.\ Koller, S.\ Hadfield, and R.\ Bowden,
``Sign language transformers: Joint end-to-end sign language recognition and translation,''
\textit{Proc.\ IEEE/CVF CVPR}, pp.\ 10023--10033, 2020.

\bibitem{koller2015continuous}
O.\ Koller, J.\ Forster, and H.\ Ney,
``Continuous sign language recognition: Towards large vocabulary statistical recognition systems handling multiple signers,''
\textit{Computer Vision and Image Understanding}, vol.\ 141, pp.\ 108--125, 2015.

\bibitem{starner1995asl}
T.\ Starner and A.\ Pentland,
``Real-time American sign language recognition from video using hidden Markov models,''
\textit{Proc.\ International Workshop on Automatic Face and Gesture Recognition}, pp.\ 265--270, 1995.

\bibitem{karpouzis2007avatar}
K.\ Karpouzis, G.\ Caridakis, S.-E.\ Fotinea, and E.\ Efthimiou,
``Educational resources and implementation of a Greek sign language synthesis architecture,''
\textit{Computers \& Education}, vol.\ 49, no.\ 1, pp.\ 54--74, 2007.

\bibitem{efthimiou2004signsynth}
E.\ Efthimiou, S.-E.\ Fotinea, T.\ Hanke, J.\ Glauert, R.\ Kennaway, and C.\ Vogler,
``The eSIGN editor: A tool for creating signed content for the Web,''
\textit{Proc.\ LREC Workshop}, 2004.

\bibitem{piao2020sign}
L.\ Piao, K.\ MacLaughlin, and M.\ Giezen,
``British Sign Language generation from text using a clip-retrieval pipeline,''
\textit{Proc.\ LREC}, pp.\ 5400--5406, 2020.

\bibitem{shardlow2014text}
M.\ Shardlow,
``A survey of automated text simplification,''
\textit{International Journal of Advanced Computer Science and Applications}, Special Issue on Natural Language Processing, pp.\ 58--70, 2014.

\bibitem{vaswani2017attention}
A.\ Vaswani, N.\ Shazeer, N.\ Parmar \textit{et al.},
``Attention is all you need,''
\textit{Advances in Neural Information Processing Systems (NeurIPS)}, vol.\ 30, 2017.

\bibitem{react2023hooks}
Meta Platforms, Inc.,
``Rules of Hooks -- React Documentation,''
\url{https://react.dev/reference/rules/rules-of-hooks}, 2023.

\bibitem{fastapi2023}
S.\ Ram\'irez,
``FastAPI: Modern, fast (high-performance) web framework for building APIs with Python 3.7+,''
\url{https://fastapi.tiangolo.com}, 2023.

\bibitem{zhang2021sign}
X.\ Zhang, C.\ Cui, and Z.\ Wang,
``A survey on sign language recognition,''
\textit{ACM Computing Surveys}, vol.\ 54, no.\ 9, pp.\ 1--37, 2022.

\bibitem{pfau2012sign}
R.\ Pfau, M.\ Steinbach, and B.\ Woll (Eds.),
\textit{Sign Language: An International Handbook}.
De Gruyter Mouton, 2012.

\end{thebibliography}

\end{document}
